\section{Implementation Details} \label{sec:implementation}

\subsection{Development Environment}

Our implementation leverages the following technology stack:

\begin{table}[h!]
\centering
\caption{Software Stack and Dependencies}
\label{tab:stack}
\begin{tabular}{|l|l|l|}
\hline
\textbf{Component} & \textbf{Library} & \textbf{Version} \\
\hline
Data Processing & Pandas, NumPy & 1.23, 1.21 \\
\hline
ML/DL Framework & scikit-learn, PyTorch & 1.2, 1.12 \\
\hline
Transformers & Hugging Face Transformers & 4.20 \\
\hline
Visualization & Matplotlib, Seaborn & 3.5, 0.12 \\
\hline
GPU Acceleration & CUDA, cuDNN & 11.6, 8.4 \\
\hline
\end{tabular}
\end{table}

All code implemented in Python 3.9+ with Jupyter Notebook for reproducible research.

\subsection{Module Structure}

The implementation is organized into modular components:

\subsubsection{Data Pipeline Module}

\begin{itemize}
  \item \textbf{Data Loading} --- Load and parse corpus from external sources
  \item \textbf{Pair Generation} --- Create document pairs with controllable plagiarism injection
  \item \textbf{Preprocessing} --- Tokenization, normalization, padding
  \item \textbf{Train/Validation/Test Splitting} --- Stratified k-fold splitting
\end{itemize}

\subsubsection{Model Implementations}

Each approach encapsulated in separate module:

\begin{enumerate}
  \item \textbf{TF-IDF Module} ($\sim$50 lines)
    \begin{itemize}
      \item Vectorizer initialization
      \item Fit on training data
      \item Inference via cosine similarity
    \end{itemize}
  
  \item \textbf{BiLSTM Module} ($\sim$200 lines)
    \begin{itemize}
      \item Vocabulary building from training corpus
      \item Model architecture definition
      \item Training loop with gradient updates
      \item Inference with probability output
    \end{itemize}
  
  \item \textbf{BERT Module} ($\sim$150 lines)
    \begin{itemize}
      \item Model initialization from Hugging Face
      \item Fine-tuning loop
      \item Tokenization and input formatting
      \item Inference with confidence scores
    \end{itemize}
\end{enumerate}

\subsubsection{Evaluation Module}

Unified metrics computation:

\begin{itemize}
  \item Accuracy, Precision, Recall, F1-Score computation
  \item ROC curve and AUC calculation
  \item Confusion matrices for error analysis
  \item Statistical significance testing (paired t-tests)
\end{itemize}

\subsubsection{Visualization Module}

\begin{itemize}
  \item Performance comparison bar charts
  \item Computational efficiency plots
  \item ROC curve plotting
  \item Confusion matrix heatmaps
  \item Radar charts for multi-metric comparison
\end{itemize}

\subsection{Hyperparameter Tuning}

Hyperparameters selected via validation set performance:

\begin{table}[h!]
\centering
\caption{Hyperparameter Exploration Results}
\label{tab:hyperparams}
\begin{tabular}{|l|c|c|c|}
\hline
\textbf{Model} & \textbf{Parameter} & \textbf{Range} & \textbf{Optimal} \\
\hline
\multirow{2}{*}{TF-IDF} & Max Features & [1000, 5000, 10000] & 5000 \\
& N-gram Range & [(1,1), (1,2), (1,3)] & (1,2) \\
\hline
\multirow{3}{*}{BiLSTM} & Hidden Dim & [32, 64, 128] & 64 \\
& Learning Rate & [0.0001, 0.001, 0.01] & 0.001 \\
& Dropout & [0.1, 0.3, 0.5] & 0.3 \\
\hline
\multirow{2}{*}{BERT} & Learning Rate & [1e-5, 2e-5, 5e-5] & 2e-5 \\
& Batch Size & [8, 16, 32] & 16 \\
\hline
\end{tabular}
\end{table}

\subsection{Reproducibility Measures}

For reproducible research:

\begin{enumerate}
  \item \textbf{Random seed} --- Fixed to 42 across NumPy, PyTorch, and Transformers
  \item \textbf{Deterministic operations} --- CUDA deterministic mode enabled
  \item \textbf{Version pinning} --- All dependencies version-locked in requirements.txt
  \item \textbf{Checkpointing} --- Model weights saved after each epoch
  \item \textbf{Logging} --- Comprehensive training/validation metrics logged
\end{enumerate}

\subsection{Computational Requirements}

\begin{table}[h!]
\centering
\caption{Training Time and Resource Requirements}
\label{tab:resources}
\begin{tabular}{|l|c|c|c|}
\hline
\textbf{Model} & \textbf{Training Time} & \textbf{GPU Memory} & \textbf{Inference Time} \\
\hline
TF-IDF & \textless 1s & -- & \textless 1ms \\
BiLSTM & 5-10 minutes & 2-3 GB & 100-200ms \\
BERT & 15-30 minutes & 6-8 GB & 300-500ms \\
\hline
\end{tabular}
\end{table}

These measurements obtained on NVIDIA GPU with 24GB memory. CPU-only inference possible but significantly slower (5-10x for BERT).
